\documentclass[11pt]{article}

\usepackage[margin=1in]{geometry}
\usepackage{amsmath,amssymb}

\begin{document}

\section*{Use of Complex Electric Fields in Linear and Nonlinear Optics}

\section{Linear Case}

The physical electric field is always real.  For a monochromatic or narrowband field, it is convenient to introduce a complex representation,
\begin{equation}
E_{\mathrm{phys}}(z,t)
=
\operatorname{Re}\!\left[\tilde{E}(z,t)\right].
\end{equation}

For example, one may write
\begin{equation}
\tilde{E}(z,t)
=
\mathcal{E}(z)e^{-i\omega t},
\end{equation}
or, for a plane wave,
\begin{equation}
\tilde{E}(z,t)
=
E_0 e^{i(kz-\omega t)}.
\end{equation}

The usefulness of this representation follows from the linearity of the differential operators that appear in Maxwell's equations.

\subsection{Linearity of Time and Spatial Derivatives}

Both time and spatial derivatives are linear operators. For constants \(a\) and \(b\),
\begin{equation}
\frac{\partial}{\partial t}(aE_1+bE_2)
=
a\frac{\partial E_1}{\partial t}
+
b\frac{\partial E_2}{\partial t},
\end{equation}
and
\begin{equation}
\frac{\partial}{\partial z}(aE_1+bE_2)
=
a\frac{\partial E_1}{\partial z}
+
b\frac{\partial E_2}{\partial z}.
\end{equation}

The same property holds for higher derivatives and for \(\nabla\), \(\nabla\cdot\), \(\nabla\times\), and \(\nabla^2\). Hence, for any linear differential operator \(\hat{L}\) with real coefficients,
\begin{equation}
\hat{L}\,
\operatorname{Re}\!\left[\tilde{E}\right]
=
\operatorname{Re}\!\left[
\hat{L}\tilde{E}
\right].
\end{equation}

Therefore, if
\begin{equation}
\hat{L}E_{\mathrm{phys}}=0,
\end{equation}
one may solve
\begin{equation}
\hat{L}\tilde{E}=0
\end{equation}
and recover
\begin{equation}
\boxed{
E_{\mathrm{phys}}
=
\operatorname{Re}\!\left[\tilde{E}\right]
}.
\end{equation}

\subsection{Why Complex Exponentials Are Convenient}

For a plane-wave complex field
\begin{equation}
\tilde{E}(z,t)
=
E_0e^{i(kz-\omega t)},
\end{equation}
spatial differentiation becomes multiplication by \(ik\),
\begin{equation}
\frac{\partial\tilde{E}}{\partial z}
=
ik\tilde{E},
\end{equation}
and
\begin{equation}
\frac{\partial^2\tilde{E}}{\partial z^2}
=
-k^2\tilde{E}.
\end{equation}

Likewise, time differentiation becomes multiplication by \(-i\omega\),
\begin{equation}
\frac{\partial\tilde{E}}{\partial t}
=
-i\omega\tilde{E},
\end{equation}
and
\begin{equation}
\frac{\partial^2\tilde{E}}{\partial t^2}
=
-\omega^2\tilde{E}.
\end{equation}

Thus the differential wave equation can be converted into an algebraic equation for \(k\) and \(\omega\).  This is one of the principal advantages of the complex representation.

\subsection{Linear Maxwell Wave Equation}

As an example, consider the linear wave equation
\begin{equation}
\nabla^2E
-
\frac{1}{c^2}
\frac{\partial^2E}{\partial t^2}
=
0.
\end{equation}

Both terms are linear in \(E\).  Substituting the complex plane wave gives
\begin{equation}
-k^2\tilde{E}
+
\frac{\omega^2}{c^2}\tilde{E}
=
0,
\end{equation}
and hence
\begin{equation}
k^2=\frac{\omega^2}{c^2}.
\end{equation}

The real physical wave is then obtained by taking the real part of the complex solution.

\subsection{Linear Propagation in a Dispersive or Absorbing Medium}

A medium can still be linear even when its susceptibility and refractive index are complex.  Linearity means that the polarization depends linearly on the field, not that the proportionality coefficient must be real.

For a monochromatic field one may write
\begin{equation}
\tilde{P}(\omega)
=
\epsilon_0\tilde{\chi}(\omega)\tilde{E}(\omega),
\end{equation}
where
\begin{equation}
\tilde{\chi}(\omega)
=
\chi'(\omega)+i\chi''(\omega).
\end{equation}

The corresponding complex refractive index is
\begin{equation}
\tilde{n}
=
n'+in'',
\end{equation}
with
\begin{equation}
\tilde{n}^{\,2}
=
1+\tilde{\chi}
\end{equation}
for a nonmagnetic medium.

The complex wave number is
\begin{equation}
\tilde{k}
=
k_0\tilde{n}
=
k_0(n'+in''),
\end{equation}
where
\begin{equation}
k_0=\frac{\omega}{c}.
\end{equation}

The complex field is therefore
\begin{equation}
\tilde{E}(z,t)
=
E_0 e^{i(\tilde{k}z-\omega t)},
\end{equation}
or
\begin{equation}
\tilde{E}(z,t)
=
E_0
e^{-k_0n''z}
e^{i(k_0n'z-\omega t)}.
\end{equation}

The measurable field is
\begin{equation}
E_{\mathrm{phys}}(z,t)
=
\operatorname{Re}\!\left[\tilde{E}(z,t)\right]
=
E_0
e^{-k_0n''z}
\cos\!\left(k_0n'z-\omega t\right).
\end{equation}

Thus \(n'\) determines the phase propagation, while \(n''\) determines attenuation.

The appearance of complex coefficients in the frequency-domain equations does not make the response nonlinear.  The system remains linear as long as superposition holds.  For a real physical medium, the frequency-domain response also satisfies the conjugate-symmetry condition
\begin{equation}
\tilde{\chi}(-\omega)
=
\tilde{\chi}^{*}(\omega),
\end{equation}
which guarantees that the reconstructed time-domain polarization is real.

\subsection{What Would Make the Equation Nonlinear?}

A derivative itself does not make an equation nonlinear.  For example,
\begin{equation}
\frac{\partial E}{\partial z}
\end{equation}
and
\begin{equation}
\frac{\partial^2 E}{\partial z^2}
\end{equation}
are linear operations.

However, products involving the field and its derivatives are generally nonlinear.  For example,
\begin{equation}
E\frac{\partial E}{\partial z}
\end{equation}
is nonlinear in \(E\), and so is
\begin{equation}
E^2.
\end{equation}

The defining criterion is therefore superposition.  An operator \(\hat{L}\) is linear if
\begin{equation}
\boxed{
\hat{L}(aE_1+bE_2)
=
a\hat{L}(E_1)
+
b\hat{L}(E_2)
}.
\end{equation}

\section{Why Extra Care Is Needed in Nonlinear Problems}

In a nonlinear problem, nonlinear operations do not generally commute with taking the real part. For example,
\begin{equation}
\left[
\operatorname{Re}\tilde{E}
\right]^2
\neq
\operatorname{Re}\!\left[
\tilde{E}^{\,2}
\right].
\end{equation}

Therefore, one cannot in general replace the real physical field by a single complex field, apply the nonlinear operation, and take the real part only at the end.

A correct and systematic representation is to separate the real field into positive- and negative-frequency parts,
\begin{equation}
E_{\mathrm{phys}}(t)
=
E^{(+)}(t)+E^{(-)}(t),
\end{equation}
where
\begin{equation}
E^{(+)}(t)
=
\frac{1}{2}\mathcal{E}e^{-i\omega t},
\end{equation}
and
\begin{equation}
E^{(-)}(t)
=
\frac{1}{2}\mathcal{E}^{*}e^{i\omega t}.
\end{equation}

Because the physical field is real,
\begin{equation}
E^{(-)}(t)
=
\left[E^{(+)}(t)\right]^*.
\end{equation}

Equivalently,
\begin{equation}
E_{\mathrm{phys}}(t)
=
\frac{1}{2}
\left[
\mathcal{E}e^{-i\omega t}
+
\mathcal{E}^{*}e^{i\omega t}
\right].
\end{equation}

\section{Example: Second-Order Nonlinearity}

Consider a second-order nonlinear polarization,
\begin{equation}
P^{(2)}(t)
=
\epsilon_0\chi^{(2)}
E_{\mathrm{phys}}^2(t).
\end{equation}

Using the full real-field representation,
\begin{equation}
E_{\mathrm{phys}}(t)
=
\frac{1}{2}
\left[
\mathcal{E}e^{-i\omega t}
+
\mathcal{E}^{*}e^{i\omega t}
\right],
\end{equation}
we obtain
\begin{equation}
E_{\mathrm{phys}}^2(t)
=
\frac{1}{4}
\left[
\mathcal{E}^2e^{-2i\omega t}
+
2|\mathcal{E}|^2
+
(\mathcal{E}^{*})^2e^{2i\omega t}
\right].
\end{equation}

Therefore,
\begin{equation}
P^{(2)}(t)
=
\frac{\epsilon_0\chi^{(2)}}{4}
\left[
\mathcal{E}^2e^{-2i\omega t}
+
2|\mathcal{E}|^2
+
(\mathcal{E}^{*})^2e^{2i\omega t}
\right].
\end{equation}

The terms proportional to \(e^{\mp 2i\omega t}\) describe second-harmonic generation, while the frequency-independent term describes optical rectification.

If one incorrectly used only
\begin{equation}
\tilde{E}(t)
=
\mathcal{E}e^{-i\omega t}
\end{equation}
and calculated
\begin{equation}
\tilde{E}^{\,2}(t)
=
\mathcal{E}^2e^{-2i\omega t},
\end{equation}
the zero-frequency contribution would be missed.

\section{Example: Third-Order Nonlinearity}

For a third-order nonlinear polarization,
\begin{equation}
P^{(3)}(t)
=
\epsilon_0\chi^{(3)}
E_{\mathrm{phys}}^3(t),
\end{equation}
the cubic field is
\begin{equation}
E_{\mathrm{phys}}^3(t)
=
\frac{1}{8}
\left[
\mathcal{E}^3e^{-3i\omega t}
+
3|\mathcal{E}|^2\mathcal{E}e^{-i\omega t}
+
3|\mathcal{E}|^2\mathcal{E}^{*}e^{i\omega t}
+
(\mathcal{E}^{*})^3e^{3i\omega t}
\right].
\end{equation}

The component at \(3\omega\) gives third-harmonic generation, while the component at the original frequency \(\omega\) is proportional to
\begin{equation}
P^{(3)}(\omega)
\propto
|\mathcal{E}|^2\mathcal{E}.
\end{equation}

This term is responsible for effects such as the optical Kerr effect and self-phase modulation.

\section{Complex Envelopes in Nonlinear Optics}

Complex notation remains very useful in nonlinear optics. A common representation is
\begin{equation}
E_{\mathrm{phys}}(z,t)
=
\frac{1}{2}
\left[
\mathcal{E}(z,t)
e^{i(kz-\omega t)}
+
\mathcal{E}^{*}(z,t)
e^{-i(kz-\omega t)}
\right],
\end{equation}
where \(\mathcal{E}(z,t)\) is a slowly varying complex envelope.

After the nonlinear polarization is expanded correctly, one can isolate the frequency component of interest. For example, in a Kerr medium the component near the carrier frequency often has the form
\begin{equation}
P^{(3)}_{\omega}
\propto
|\mathcal{E}|^2\mathcal{E}
e^{i(kz-\omega t)}.
\end{equation}

The corresponding envelope equation can then contain a term such as
\begin{equation}
\frac{\partial \mathcal{E}}{\partial z}
\propto
i|\mathcal{E}|^2\mathcal{E}.
\end{equation}

Thus complex amplitudes remain powerful in nonlinear problems, provided the nonlinear mixing of positive- and negative-frequency components is handled first.

\section{Multiple Frequencies}

For a field containing several frequencies, write
\begin{equation}
E_{\mathrm{phys}}(t)
=
\frac{1}{2}
\sum_j
\left[
\mathcal{E}_j e^{-i\omega_j t}
+
\mathcal{E}_j^{*}e^{i\omega_j t}
\right].
\end{equation}

A second-order nonlinearity generates frequency components satisfying
\begin{equation}
\omega_3=\omega_1+\omega_2.
\end{equation}

Because negative frequencies are also included,
\begin{equation}
E(-\omega)=E^{*}(\omega),
\end{equation}
the same formalism automatically describes both sum- and difference-frequency generation.

\section{Complex Amplitudes in Linear and Nonlinear Optics}

Complex amplitudes or complex envelopes are useful in both linear and nonlinear optics.  For a field containing several frequency components, one may write
\begin{equation}
E_{\mathrm{phys}}(z,t)
=
\frac{1}{2}
\sum_j
\left[
\mathcal{E}_j(z,t)e^{i(k_j z-\omega_j t)}
+
\mathcal{E}_j^{*}(z,t)e^{-i(k_j z-\omega_j t)}
\right],
\end{equation}
where each \(\mathcal{E}_j(z,t)\) is a complex amplitude or envelope containing the amplitude and phase of the corresponding frequency component.

In a linear medium, each frequency component evolves independently.  For example,
\begin{equation}
\tilde{P}(\omega_j)
=
\epsilon_0
\tilde{\chi}^{(1)}(\omega_j)
\tilde{E}(\omega_j).
\end{equation}

Thus a field component at frequency \(\omega_j\) produces a polarization at the same frequency \(\omega_j\).  One may therefore solve separately for the complex amplitude \(\mathcal{E}_j\) of each component.

In a nonlinear medium, the different frequency components can mix.  For example, a second-order polarization has the form
\begin{equation}
P^{(2)}
=
\epsilon_0\chi^{(2)}
E_{\mathrm{phys}}^2,
\end{equation}
and can generate components at frequencies such as
\begin{equation}
2\omega_1,\qquad
2\omega_2,\qquad
\omega_1+\omega_2,\qquad
\omega_1-\omega_2.
\end{equation}

The nonlinear polarization should therefore first be expanded using the full real field,
\begin{equation}
E_{\mathrm{phys}}
=
E^{(+)}+E^{(-)},
\qquad
E^{(-)}
=
\left(E^{(+)}\right)^*.
\end{equation}

After the generated frequency components have been identified, each relevant component can again be represented by its own complex amplitude or complex envelope.

Thus the use of complex amplitudes is not unique to nonlinear optics.  The essential distinction is
\begin{equation}
\boxed{
\text{linear medium: frequencies evolve independently}
}
\end{equation}
whereas
\begin{equation}
\boxed{
\text{nonlinear medium: frequencies can mix and generate new frequencies}.
}
\end{equation}

\section{Comparison of Linear and Nonlinear Cases}

\begin{table}[ht]
\centering
\renewcommand{\arraystretch}{1.35}
\begin{tabular}{p{0.28\textwidth} p{0.30\textwidth} p{0.30\textwidth}}
\hline
\textbf{Feature} & \textbf{Linear case} & \textbf{Nonlinear case} \\
\hline
Complex amplitudes or envelopes & Used for each frequency component & Used for each input and generated frequency component \\
\hline
Material response & Polarization is linear in the field & Polarization contains powers or products of fields \\
\hline
Typical relation & \(P^{(1)}(\omega)=\epsilon_0\chi^{(1)}(\omega)E(\omega)\) & \(P^{(n)}\propto E^n\) \\
\hline
Frequency mixing & No & Yes \\
\hline
New frequencies generated & No & Yes, in general \\
\hline
Treatment of different frequencies & Each frequency can be solved independently & Frequencies are coupled through the nonlinear polarization \\
\hline
Need to retain \(E^{(+)}\) and \(E^{(-)}\) explicitly & Usually unnecessary for a single linear frequency component & Generally necessary before forming nonlinear products \\
\hline
Final representation & Complex amplitude or envelope may remain the useful final result & Complex amplitudes or envelopes of selected generated components are usually the useful final results \\
\hline
Physical electric field & Obtained by reconstructing the real field when needed & Obtained by reconstructing the real field when needed \\
\hline
Typical observable & \( |\mathcal{E}|^2 \), phase, transmission & \( |\mathcal{E}_j|^2 \), conversion efficiency, phase, nonlinear phase shift \\
\hline
\end{tabular}
\caption{Comparison of the use of complex representations in linear and nonlinear optics.}
\end{table}

\section{Practical Procedure}

For a linear problem:

\begin{enumerate}
    \item decompose the real field into the frequency components of interest;
    \item represent each component by a complex amplitude or envelope;
    \item solve the linear propagation equation independently for each frequency;
    \item reconstruct the real field only if the instantaneous physical field is required.
\end{enumerate}

For a nonlinear problem:

\begin{enumerate}
    \item write the full real field in terms of positive- and negative-frequency components;
    \item substitute it into the nonlinear polarization or nonlinear equation;
    \item expand the nonlinear products;
    \item group the resulting terms by frequency;
    \item identify the generated frequency components relevant to the process;
    \item represent each selected component by its own complex amplitude or envelope and solve the corresponding coupled propagation equations;
    \item reconstruct the real field only if the instantaneous physical field is required.
\end{enumerate}

In both cases, the complex amplitude or envelope may remain complex throughout the calculation.  Physical observables are real quantities obtained either by reconstructing the real field,
\begin{equation}
E_{\mathrm{phys}}
=
E^{(+)}+E^{(-)},
\end{equation}
or directly from real combinations such as
\begin{equation}
I
\propto
|\mathcal{E}|^2.
\end{equation}


\end{document}
